%%%%%%%%%%%%%%%%%%%%%%%%%%%%%%%%%%%%%%%%%%%%%%%%%%%%%%%%%%%%%%
%% Rozdział 8: Wnioski
%%%%%%%%%%%%%%%%%%%%%%%%%%%%%%%%%%%%%%%%%%%%%%%%%%%%%%%%%%%%%%

\section{Wnioski}
\label{chap:conclusions}

\subsection{Realizacja celów pracy}
\label{sec:conclusions-goals}

W podrozdziale~\ref{sec:goals-scope} określono pięć celów projektu. Dwa z nich
zostały w pełni zrealizowane w ramach prac przedstawionych w niniejszej pracy,
a trzy pozostają niezrealizowane w chwili jej pisania.

Cel dotyczący platformy do anotacji został osiągnięty dzięki systemowi opisanemu
w rozdziale~\ref{chap:annotation-platform}: działającej, zdecentralizowanej
aplikacji do gromadzenia danych, która obsługuje zarządzanie zbiorami danych,
nagraniami wideo, klasami i typami gestów, asynchroniczną ekstrakcję punktów
charakterystycznych z wykorzystaniem MediaPipe oraz przeglądarkowy edytor
anotacji klipów. Całość można uruchomić lokalnie za pomocą Docker Compose.

Cel dotyczący rozpoznawania izolowanych znaków w procesie wideo--glosy został
osiągnięty dzięki systemowi opisanemu
w rozdziale~\ref{chap:isolated-gloss-recognition}: wielocelowej architekturze
opartej na wyszukiwaniu, ocenionej z wykorzystaniem pięciu enkoderów czasowych
na podzbiorze ASL Citizen obejmującym 32 glosy. Enkoder Transformer osiągnął
dokładność 1-NN wynoszącą 90,75\% oraz czułość dla każdej klasy przekraczającą
0,75 dla wszystkich ocenianych glos.

\todo{DO UZUPEŁNIENIA: Po ukończeniu rozdziałów~\ref{chap:text-translation},
\ref{chap:sentence-level-translation} i~\ref{chap:educational-platform} należy
podsumować, w jaki sposób i w jakim stopniu osiągnięto każdy z pozostałych celów
-- tłumaczenie między glosami a tekstem w obu kierunkach, całościowe tłumaczenie
na poziomie zdań oraz stworzenie edukacyjnej platformy demonstracyjnej --
z odwołaniem do odpowiednich podrozdziałów prezentujących wyniki.}

\subsection{Ograniczenia}
\label{sec:conclusions-limitations}

Głównym ograniczeniem systemu rozpoznawania izolowanych glos, zidentyfikowanym
na podstawie analizy macierzy pomyłek, czułości dla poszczególnych klas,
wizualizacji t-SNE oraz śledzenia odległości w parach pozytywnych i negatywnych
(podrozdział~\ref{sec:gloss-recognition-discussion}), jest zmienność
wewnątrzklasowa wynikająca z różnic w wykonywaniu znaków przez poszczególne osoby
migające. Osoby niewidziane podczas trenowania generują bardziej rozproszone
osadzenia niż osoby widziane, mimo że ogólna struktura klastrów zachowuje zdolność
uogólniania. Obecne porównanie modeli jest ponadto ograniczone do podzbioru
ASL Citizen obejmującego 32 glosy oraz do izolowanych, uprzednio wyodrębnionych
znaków, a nie ciągłych wypowiedzi migowych.

Do znanych ograniczeń platformy do anotacji
(podrozdział~\ref{sec:annotation-system-architecture}) należą nieukończona obsługa
przetwarzania eksportu zbiorów danych, metadane nagrań (FPS, czas trwania), które
nie są jeszcze w pełni wyznaczane na podstawie przesłanych plików, brak
uwierzytelniania i autoryzacji, ograniczona walidacja zakresów klipów i spójności
zbiorów danych oraz ograniczony zakres testów automatycznych.

\todo{DO UZUPEŁNIENIA: Po ukończeniu rozdziałów~\ref{chap:text-translation},
\ref{chap:sentence-level-translation} i~\ref{chap:educational-platform} należy
dodać właściwe im ograniczenia (np.\ wielkość i zakres korpusu glos i tekstu,
propagację błędów między etapami rozpoznawania i tłumaczenia w całościowym potoku
przetwarzania, ograniczenia użyteczności platformy edukacyjnej).}

\subsection{Kierunki dalszych prac}
\label{sec:conclusions-future-work}

W podrozdziale~\ref{sec:gloss-recognition-discussion} wskazano kilka konkretnych
rozszerzeń systemu rozpoznawania izolowanych glos: zastąpienie lub uzupełnienie
klasyfikatora KNN dopasowaniem sekwencji opartym na dynamicznym dopasowaniu czasu
(Dynamic Time Warping, DTW) na etapie wyszukiwania; zastąpienie funkcji straty
rekonstrukcji MSE metodą soft-DTW jako różniczkowalną funkcją celu uczenia;
zastosowanie FAISS do skalowalnego, przybliżonego wyszukiwania najbliższych
sąsiadów poza obecnym benchmarkiem obejmującym 32 glosy; wprowadzenie grafowych
sieci splotowych w celu jawnego kodowania topologii grafu szkieletu zamiast
polegania na spłaszczonych wektorach punktów charakterystycznych; ocenę sieci
LSTM operujących na kwaternionach do kodowania zależności czasowych
z uwzględnieniem rotacji; oraz rozszerzenie systemu z rozpoznawania izolowanego
na ciągłe za pomocą CTC, stanowiące podstawę potoku przetwarzania na poziomie
zdań opisanego w rozdziale~\ref{chap:sentence-level-translation}. Warto zaznaczyć,
że od czasu sporządzenia raportu badawczego stanowiącego podstawę niniejszej pracy
w repozytorium projektu zaimplementowano już demonstrację rozpoznawania na żywo
z kamery internetowej, łączącą enkoder Transformer z wyszukiwaniem za pomocą
FAISS, zanim część powyższych rozszerzeń została formalnie opisana.

W przypadku platformy do anotacji dalsze prace powinny obejmować ukończenie
potoku przetwarzania eksportu zbiorów danych, usprawnienie śledzenia statusów
zadań importu i eksportu, wyznaczanie metadanych nagrań bezpośrednio na podstawie
przesłanych plików zamiast stosowania stałych wartości, dodanie uwierzytelniania
 i autoryzacji oraz rozszerzenie zakresu testów automatycznych obejmujących API,
proces roboczy i edytor we frontendzie.

\todo{DO UZUPEŁNIENIA: Po ukończeniu rozdziałów~\ref{chap:text-translation},
\ref{chap:sentence-level-translation} i~\ref{chap:educational-platform} należy
dodać właściwe im kierunki dalszych prac.}